\documentclass[12pt,a4paper]{article}
\usepackage[T1]{fontenc}
\usepackage[utf8]{inputenc}
\usepackage[polish]{babel}
\usepackage[a4paper,margin=2.5cm]{geometry}
\usepackage{enumitem}
\usepackage{hyperref}

\title{Raport postępów projektu\\System zarządzania domową biblioteką}
\author{Imię i nazwisko}
\date{\today}

\begin{document}
\maketitle

\section{Cel projektu}
Celem projektu jest stworzenie konsolowej aplikacji w języku C++ do zarządzania domową biblioteką.
Aplikacja ma wykorzystywać podejście obiektowe (OOP), umożliwiać pełny zakres operacji CRUD
na danych oraz dostarczać podstawowe statystyki i trwałość danych w plikach CSV.

\section{Aktualnie zaimplementowane funkcjonalności}
\begin{itemize}[leftmargin=*]
    \item Obsługa encji: autorzy, kategorie i pozycje biblioteczne.
    \item Pełne CRUD dla autorów, kategorii i pozycji.
    \item Wyszukiwanie pozycji po tytule, autorze lub kategorii.
    \item Statystyki: liczba pozycji/autorów/kategorii, średnia ocen, zestawienia per kategoria.
    \item Interfejs konsolowy z menu i podmenu do zarządzania danymi.
\end{itemize}

\section{Architektura i klasy (OOP)}
Projekt ma strukturę warstwową:
\begin{itemize}[leftmargin=*]
    \item \texttt{src/models}: klasy modelu domenowego (\texttt{Author}, \texttt{Category}, \texttt{Item}).
    \item \texttt{src/core}: logika biznesowa i operacje na kolekcjach (\texttt{Library}).
    \item \texttt{src/utils}: obsługa zapisu/odczytu danych (\texttt{FileStorage}).
    \item \texttt{src/main.cpp}: warstwa UI (interfejs tekstowy i nawigacja po menu).
\end{itemize}
Stosowane są podstawowe założenia OOP: enkapsulacja, podział odpowiedzialności i separacja warstw.
W obecnej wersji projektu najpełniej wykorzystywana jest enkapsulacja oraz kompozycja klas.

\section{Trwałość danych (pliki CSV)}
Dane aplikacji są utrwalane w katalogu \texttt{data/}:
\begin{itemize}[leftmargin=*]
    \item \texttt{authors.csv} -- autorzy,
    \item \texttt{categories.csv} -- kategorie,
    \item \texttt{items.csv} -- pozycje biblioteczne.
\end{itemize}
Aplikacja obsługuje automatyczne i ręczne wczytywanie/zapisywanie danych,
dzięki czemu stan biblioteki jest zachowany pomiędzy uruchomieniami.

\section{Instrukcja budowania i uruchamiania}
\subsection*{Linux/macOS}
\begin{verbatim}
cmake -B build -DCMAKE_BUILD_TYPE=Release
cmake --build build
./build/HomeLibrary
\end{verbatim}

\subsection*{Windows}
\begin{verbatim}
cmake -B build
cmake --build build --config Release
build\Release\HomeLibrary.exe
\end{verbatim}

\section{Kolejne kroki}
\begin{enumerate}[leftmargin=*]
    \item Rozbudowa walidacji danych wejściowych i obsługi błędów użytkownika.
    \item Dodanie testów jednostkowych kluczowych klas (modele, \texttt{Library}, \texttt{FileStorage}).
    \item Dalsze uporządkowanie warstwy UI i poprawa ergonomii menu.
    \item Ewentualna rozbudowa o funkcjonalności rozszerzone (np. GUI lub baza danych).
\end{enumerate}

\end{document}
